\chapter{绪论}

\section{研究背景}

\subsection{能源转型与海洋资源开发的共同要求}

《中华人民共和国能源法》将风能、太阳能、海洋能、氢能等纳入能源范畴，并要求加快
构建清洁低碳、安全高效的新型能源体系\cite{energy-law}。随着近岸优质场址逐步开发，
海上可再生能源利用正由单一近海风电场向深水、远岸、多能融合和综合开发延伸。深远海
资源潜力并不会自动转化为可交付价值：远岸距离、复杂海况、弱网运行以及产品交付共同
决定项目能否形成稳定收益。
\sourcetype{政府官网｜《中华人民共和国能源法》；本段只说明战略方向，不使用未经核验的资源量。}

国家能源局已将海上能源岛列入可再生能源试点示范方向，提出结合海上风电开发，融合
区域储能、制氢、海水淡化和海洋养殖等需求，探索能源资源供给与转换一体化设施
\cite{nea-demo}。2025年公开答复进一步将海上能源岛概括为依托岛屿或海上平台，集合
多种海上可再生能源、储能、氢氨醇制备、电力送出和综合服务的能源枢纽，并明确指出
国内外相关路线仍处于起步探索阶段\cite{nea-energy-island}。这说明“综合化”已有政策
场景，但完整工程方案和产业模式仍需验证。
\sourcetype{政府官网｜国家能源局可再生能源试点示范通知及海上能源岛公开答复。}

\subsection{研究对象由“发电项目”转向“价值转化系统”}

传统海上风电项目主要回答“如何发电并送出”，而深远海能源枢纽还必须回答“有限的
清洁电力在何时、以何种形态交付给谁”。本项目据此提出“源—储—算—用”四位一体
方案：以漂浮式风电为主体，配置漂浮式光伏并为波浪能、潮流能预留接口；以构网型
电池储能提供快速稳定支撑，以制氢及储氢链承担较长时间尺度的能量转移；将海底绿色
算力作为具有服务约束的柔性负荷；最终形成电力、氢氨、算力服务和海洋综合用能等
多形态输出。

其中，漂浮式风电占供能容量不低于85\%为项目当前设计方向\assumptiontag，具体容量和
占比不得在绪论中提前固化，而应由第4章容量模型和第5章场景结果回填。海底算力
PUE低于1.1同样属于项目目标\assumptiontag，应通过设备方案、海水温度条件和实测边界
验证，不能作为既成性能写入结论。
\sourcetype{项目内部框架；数值目标尚需设备商、算力运营商及场址数据校准。}

\subsection{国家能源集团的场景牵引条件}

国家能源集团已开展漂浮式风电与深海养殖融合实践，“国能共享号”验证了海上发电与
海洋空间复用的工程组织能力\cite{chnenergy-sharing}；集团在新能源、储能和可再生氢
方面亦形成了可供本项目调用的板块基础。上述事实能够证明集团具备组织跨专业示范的
条件，但不能据此推导“源储算用”完整系统已成熟，也不能替代本项目的系统验证。
\sourcetype{中央企业官网｜国家能源集团“国能共享号”公开资料。}

\section{相关方向进展与研究缺口}

\subsection{单项技术已形成不同程度的工程基础}

漂浮式风电、海上光伏、构网型储能、电解水制氢、海底数据中心、柔性直流输电和智能
调度均已有研究、示范或商业应用，但成熟度并不一致。漂浮式风电全球运行规模仍明显
小于项目储备规模，呈现“需求强、在运经验有限”的产业化早期特征\cite{irena-floating}；
构网型控制和海上制氢的成熟设备也需要在深远海弱网、高盐雾、摇摆和极端天气环境中
重新验证。具体装备证据与成熟度由第3章统一评估，本章不再罗列设备参数和案例清单。
\sourcetype{国际政府间组织报告、国际标准、核心期刊论文；详见第3章。}

\subsection{现有项目多验证单链路，完整协同仍缺证据}

当前公开项目大体可归为电网汇集、风电制氢、海洋综合利用、海底算力等单链路或双链路
模式。它们证明部分接口可行，却不能直接回答多端口共同运行时的功率竞争、履约冲突、
设备退化、储运约束与风险分配。尤其当电力、氢能和算力服务对应不同市场周期时，固定
比例分配可能在某些时段可行、在另一些时段失效。因此，本研究缺口不在于再提出一张
设备拼接图，而在于建立可执行、可比较、可回溯的统一调度与价值核算体系。

\subsection{从“技术可行”到“产业可行”缺少闭环}

新能源上网电量市场化进一步强化了价格、合同和灵活性管理的重要性\cite{ndrc-price}。
但多产品收益不能简单相加：同一单位电量只能进入一个最终交付路径，算力和制氢收益
必须扣除用电机会成本与履约成本，环境权益也必须遵守唯一性要求。现有研究常在工程、
调度和商业分析之间割裂，导致技术参数无法进入现金流，或商业假设无法回到逐时物理
约束。本项目据此把第3章的设备边界、第4章的模型、第5章的对比验证和第6章的投资
闸门连接为同一证据链。
\sourcetype{政府官网｜新能源上网电价市场化政策；项目内部研究归纳。}

\section{核心问题与现实挑战}

\subsection{远距离交付成本与单一电力收入不匹配}

深远海项目的海缆、换流、安装和运维边界随离岸距离、容量、电压等级、海床条件及海况
显著变化。初稿中“超过100 km后海缆超过1500万元/公里、损耗超过8\%”等数字尚缺
统一工程口径，现阶段不作为本报告事实结论，统一列为\assumptiontag，待目标海域的
海缆方案、工程量清单和设备报价形成后在第5、6章校准。可确认的问题是：若只有电力
外送一个出口，送出受限和市场价格波动将同时压缩资源价值，系统缺少主动切换路径。
\sourcetype{项目初始痛点与内部研究判断；具体造价、损耗和距离阈值待设计院及海缆企业一手数据校准。}

\subsection{弱网稳定需求与能量调度模型不在同一时间尺度}

深远海能源岛在弱网或孤岛模式下不能依赖陆上强电网提供电压、频率参考和故障支撑。
构网型电池储能需要承担快速功率平衡、模式切换与黑启动等职责，而制氢、算力和产品
外输主要参与分钟至小时级能量分配。小时级优化只能给出可行工作点，不能替代正序动态、
电磁暂态、硬件在环和现场试验。若混淆两个层级，将出现“能量守恒但动态不稳定”的
伪可行方案。

\subsection{绿色算力的高可靠服务与可再生能源波动存在张力}

算力并非可任意开关的普通可调负荷。训练、推理、存储与安全任务具有不同中断成本、
迁移窗口和服务等级协议。调度系统必须先满足关键任务和舱体辅助系统的供电边界，再在
可迁移、可延迟任务中释放柔性。算力负荷若被过度简化为无成本弃载，会高估协同收益；
若完全刚性化，又会低估算随电动的价值。第3章定义任务分级和接口，第4章建立约束，
第5章以服务违约指标检验。

\subsection{多产品比较需要统一计量与因果信息边界}

电力、氢能、算力和海洋用能的计量单位、合同周期与交付风险不同。模型必须把它们映射
到同一能量边界和价值口径，同时禁止使用未来已知信息指导当前决策。固定策略与优化
策略也不应混为同一比较层级：前者用于解释结构差异，后者用于评估在因果预测和滚动
更新下的可实现增益。该问题构成第4、5章的核心方法约束。

\begin{table}[H]
\centering
\caption{项目核心问题、责任章节与验证输出}
\label{tab:c1-question-map}
\begin{tabularx}{\textwidth}{p{2.2cm} X p{2.5cm} p{3.2cm}}
\toprule
核心问题 & 需要回答的内容 & 主责任章节 & 可接受证据\\
\midrule
系统如何组成 & 设备职责、接口、安全与控制边界 & 第3章 & 标准、公开工程证据、企业接口数据\\
调度如何决策 & 变量、约束、目标、信息集与滚动机制 & 第4章 & 来源明确的公式、项目核算定义、代码实现\\
方案提升多少 & 正常与极端场景下的经济、消纳、可靠和环境指标 & 第5章 & 48 h典型窗口与年度因果滚动仿真\\
项目如何落地 & 主体、合同、政策、闸门、示范与复制条件 & 第6章 & 政策文件、报价合同、企业访谈和仿真回填\\
\bottomrule
\end{tabularx}
\sourcetype{项目内部研究框架。}
\end{table}

\section{研究思路、方法与创新}

\subsection{研究问题}

本研究围绕四个递进问题展开：
\begin{enumerate}
  \item 如何定义适用于中国深远海环境的“源储算用”系统边界，并识别系统级成熟度短板；
  \item 如何在不预知未来总发电量的条件下，以逐时可用信息协调电力、储能、制氢、算力和外输；
  \item 相比纯电、纯氢和纯算三类资产基准，协同调度在经济性、消纳率、可靠性与环保性上提升多少，边界在哪里；
  \item 国家能源集团应以何种投资主体、合同组合、政策机制和阶段闸门完成示范及复制。
\end{enumerate}

\subsection{技术路线}

研究首先用政策和产业证据界定未来产业属性，再由系统架构冻结设备责任与接口；随后将
物理边界转写为多目标优化模型，建立不提前获知全年信息的因果滚动调度；第5章分别在
典型48 h、完整年度和极端事件中对基准及固定策略进行比较，并独立报告优化策略结果；
最后把性能指标回填投资、合同和政策闸门。由此形成“证据—架构—模型—验证—落地”的
闭环，而不是从设备参数直接跳到投资结论。

\subsection{拟形成的创新贡献}

\begin{enumerate}
  \item \textbf{系统组织创新：}把构网型混合储能、海底绿色算力和多形态产品交付纳入同一能源枢纽边界，并明确各模块的控制与责任接口；
  \item \textbf{调度方法创新：}在统一公共母线和能量守恒下，区分实时安全层、滚动调度层与容量规划层，约束优化仅使用决策时点可获得的信息；
  \item \textbf{评价口径创新：}用同一能量账簿比较电力、氢能、算力和海洋用能，防止储能充放电、产品收益和环境权益重复计算；
  \item \textbf{产业转化创新：}将TRL、仿真指标、长期合同和融资条件转化为阶段闸门，使示范项目具有可停止、可调整和可扩容的决策机制。
\end{enumerate}

以上为本项目拟验证的贡献，不在第1章提前宣告量化优势；最终提升幅度以第5章结果、
落地建议以第6章证据为准。

\subsection{章节安排}

第2章判断深远海能源岛是否属于值得布局的未来产业；第3章给出系统架构、关键技术和
TRL；第4章建立源储算用协同模型；第5章完成基准、固定策略、典型窗口、年度滚动及
极端事件验证；第6章形成投资、政策和产业培育路径；第7章汇总经验证结论、创新与局限。

\section*{本章小结}
\addcontentsline{toc}{section}{本章小结}

本章将研究焦点从“深远海能否发电”收敛为“清洁能源如何在安全边界内被转换为可交付、
可计量和可融资的多形态价值”。绪论只保留战略背景、问题链与研究方法；产业统计、设备
参数、公式、仿真结果和投资条款分别交由第2、3、4、5、6章承担，从源头避免重复论证。
